\subsection{赋值的扩张与完备化}

定义 3. 设 K 是域，v 是 K 上的赋值，L 是 K 的扩张。L 上的一个族 \((v_{t}^{\prime})_{t\in I}\) （由扩张 v 的赋值组成），使 L 上每个扩张 v 的赋值都等价于唯一的 \(v_{t}^{\prime}\) ，称为 v 到 L 的扩张的完全系。

命题 2. 设 K 是域，v 是 K 上的赋值，\(\hat{K}\) 是 K 关于 v 的完备化，\(\hat{v}\) 是 v 到 \(\hat{K}\) 的连续延拓，L 是 K 的次数为 n 的有限扩张。

\begin{enumerate}
\def\labelenumi{(\alph{enumi})}
\tightlist
\item
  设 \(v'\) 是 \(\mathbf{L}\) 上扩张 \(v\) 的赋值；用 \(\hat{\mathbf{L}}_{v'}\) 表示 \(\mathbf{L}\) 关于 \(v'\) 的完备化，\(\hat{v}'\) 是 \(v'\) 到 \(\hat{\mathbf{L}}_{v'}\) 的连续延拓；把 \(\hat{\mathbf{K}}\) 与 \(\mathbf{K}\) 在 \(\hat{\mathbf{L}}_{v'}\) 中的闭包视为同一，
\end{enumerate}

\begin{enumerate}
\def\labelenumi{(\arabic{enumi})}
\setcounter{enumi}{3}
\tightlist
\item
\end{enumerate}

\[
e (\hat {v} ^ {\prime} / \hat {v}) = e (v ^ {\prime} / v), \quad f (\hat {v} ^ {\prime} / \hat {v}) = f (v ^ {\prime} / v),\tag{5}
\]

\[
[ \hat {\mathbf {L}} _ {v ^ {\prime}}: \hat {\mathbf {K}} ] \leqslant n,\tag{6}
\]

\[
e (v ^ {\prime} / v) f (v ^ {\prime} / v) \leqslant [ \hat {\mathbf {L}} _ {v ^ {\prime}}: \hat {\mathbf {K}} ].
\]

\begin{enumerate}
\def\labelenumi{(\alph{enumi})}
\setcounter{enumi}{1}
\tightlist
\item
  \(\mathbf{L}\) 上扩张一个不是非真的赋值 \(v\) 的两两独立赋值组成的每个集合都有限。设 \(v_{1}',\ldots ,v_{s}'\) 表示 \(\mathbf{L}\) 上扩张 v 且使 L 上每个扩张 v 的赋值都相依于某个 \(v_{i}^{\prime}\) 的两两独立赋值；设 \(L_{i}\) 是赋以由 \(v_{i}^{\prime}\) 定义的拓扑的域 L ，\(\hat{L}_{i}\) 是其完备化；记 \(n_{i} = [\hat{L}_{i}: \hat{K}]\) 。那么典范映射
\end{enumerate}

\[
\phi : \hat {\mathbf {K}} \otimes_ {\mathbf {K}} \mathbf {L} \rightarrow \prod_ {i = 1} ^ {s} \hat {\mathbf {L}} _ {i}
\]

（把对角映射 \(\mathbf{L} \to \prod_{i=1}^{s} \mathbf{L}_i\) 连续延拓）是满射，其核是 \(\hat{\mathbf{K}} \otimes_{\mathbf{K}} \mathbf{L}\) 的 Jacobson 根，且

\[
\sum_ {i = 1} ^ {s} n _ {i} \leqslant n.\tag{7}
\]

先证明 (a)。设 \(v\) 不是非真的。由于 \(v\) 与 \(\hat{v}\) （分别 \(v'\) 与 \(\hat{v}'\) ）有相同的阶群与相同的剩余域（§ 5 第 3 节命题 5 (b) 与 (f)），(4) 成立。由它借助引理 2 推出 (6)。最后，向量子 \(\hat{K}\) -空间——\(\hat{L}_{v'}\) 的、由 \(L\) 生成的——是闭的（§ 5 第 2 节命题 4 的推论）且处处稠密，因而等于 \(\hat{L}_{v'}\) ；这证明了 (5)。

现在转到 (b)。仍可设 \(v\) 不是非真的。设 \((v_1', \ldots, v_r')\) 是 \(L\) 上扩张 \(v\) 的两两独立赋值的任意有限族。\(L\) 在对角映射下在 \(\prod_{i=1}^{r} L_i\) 中的像处处稠密（§ 7 第 2 节定理 1），且 \(\prod_{i=1}^{r} L_i\) 在 \(\prod_{i=1}^{r} \hat{L}_i\) 中稠密。因而 \(\hat{K} \otimes_K L\) 在 \(\prod_{i=1}^{r} \hat{L}_i\) 中的典范像处处稠密。另一方面，这个像是一个向量子 \(\hat{K}\) -空间（\(\prod_{i=1}^{r} \hat{L}_i\) 的）；由 (5)，\(\prod_{i=1}^{r} \hat{L}_i\) 在 \(\hat{K}\) 上有限维，故 \(\hat{K} \otimes_K L\) 的像是闭的（§ 5 第 2 节命题 4 的推论），因而等于 \(\prod_{i=1}^{r} \hat{L}_i\) 。由于 \(\hat{K} \otimes_K L\) 在 \(\hat{K}\) 上的维数是 \(n, \sum_{i=1}^{r} n_i \leqslant n\) 。这特别表明整数 \(r\) 以 \(n\) 为上界，并证明了 (b) 的第一个断言。

现在取陈述中的 \((v_{1}',\ldots ,v_{s}^{\prime})\) 。

\[
\phi : \hat {\mathbf {K}} \otimes_ {\mathbf {K}} \mathbf {L} \rightarrow \prod_ {i = 1} ^ {s} \hat {\mathbf {L}} _ {i}
\]

为满射这一事实以及关系 (7) 都已证明。剩下的要验证 \(\mathfrak{n}\) （\(\phi\) 的核）是 Jacobson 根 \(\mathfrak{r}\) （\(\hat{\mathbf{K}}\otimes_{\mathbb{K}}\mathbf{L}\) 的）。由于 \(\prod_{i = 1}\hat{\mathbf{L}}_i\) 半单，有 \(\mathfrak{r}\subset \mathfrak{n}\) 。另一方面，对每个极大理想 \(\mathfrak{m}\) （\(\hat{\mathbf{K}}\otimes_{\mathbb{K}}\mathbf{L}\) 的），商域 \(\mathbf{L}(\mathfrak{m}) = (\hat{\mathbf{K}}\otimes_{\mathbb{K}}\mathbf{L}) / \mathfrak{m}\) 是 \(\hat{\mathbf{K}}\) 与 \(\mathbf{L}\) 在 \(\mathbf{K}\) 上的合成扩张（《代数》，

第八章 § 8 命题 1）。存在一个赋值 \(w\) （在 \(\mathbf{L}(\mathfrak{m})\) 上扩张 \(\hat{v}\) ）（§ 3 第 3 节命题 5）；\(v'\) （\(w\) 在 \(\mathbf{L}\) 上的限制）扩张 \(v\) 。由于 \([\mathbf{L}(\mathfrak{m}): \hat{\mathbf{K}}]\) 有限，\(\mathbf{L}(\mathfrak{m})\) 关于 \(w\) 完备（§ 5 第 2 节命题 4）。而现在 \(\mathbf{L}\) 在 \(\mathbf{L}(\mathfrak{m})\) 中的闭包是包含 \(\hat{\mathbf{K}}\) 与 \(\mathbf{L}\) 的域，因而等于 \(\mathbf{L}(\mathfrak{m})\) 。于是 \(\mathbf{L}(\mathfrak{m})\) 与完备化 \(\hat{\mathbf{L}}_{v'}\) 视为同一，且 \(m\) 是 \(\hat{\mathbf{K}} \otimes_{\mathbb{K}} \mathbf{L}\) 到 \(\hat{\mathbf{L}}_{v'}\) 上的典范映射的核。现在，由假设，存在指标 \(i\) 使得 \(v'\) 与 \(v_i'\) 相依；于是 \(\mathbf{L}_{v'} = \mathbf{L}_i\) （§ 7 第 2 节命题 3）。这样 \(n \subset m\) ，这证明了 \(n \subset r\) 并完成了证明。

推论 1. 若 K 关于 v 完备且 v 不是非真的，则 L 上扩张 v 的两个赋值相依。

这由 \(\hat{\mathbf{K}}\otimes_{\mathbb{K}}\mathbf{L} = \mathbf{L}\) 得出。

推论 2. 若 \(\hat{\mathbf{K}}\) 或 \(\mathbf{L}\) 在 \(\mathbf{K}\) 上可分，则典范映射

\[
\phi : \hat {\mathbf {K}} \otimes_ {\mathbf {K}} \mathbf {L} \rightarrow \prod_ {i = 1} ^ {s} \hat {\mathbf {L}} _ {i}
\]

是同构。

这时 \(\hat{\mathbf{K}}\otimes_{\mathbf{K}}\mathbf{L}\) 的 Jacobson 根为零（《代数》第八章 § 7 第 3 节定理 1）。

注. 命题 2 表明，\(\hat{\mathbf{K}}\) 与 \(\mathbf{L}\) 在 \(\mathbf{K}\) 上的每个合成扩张（《代数》第八章 § 8）同构于某个完备化 \(\hat{\mathbf{L}}_t\) ，且这些完备化是两两不同构的合成扩张。
% semantic-fragments: legacy-input-seam
\relax{}
